% One Britannica scan as extracted by Nemotron Parse 3.5. Filled by src/nemotron_challenge/stages/latex.py.
% The TITLE, BODY and SOURCE markers (each wrapped in two percent signs) are replaced. Compile with XeTeX (tectonic).
\documentclass[10pt]{article}
\usepackage[a4paper,margin=18mm]{geometry}
\usepackage{fontspec}
\usepackage{multicol}
\usepackage{graphicx}
\usepackage[hidelinks]{hyperref}
\setlength{\parindent}{0pt}
\setlength{\parskip}{4pt}
\setlength{\columnsep}{8mm}
\begin{document}
\section*{Historical source: ATARNEA}
\begin{multicols}{2}
\subsection*{\textit{continued from the previous page}}
Ataraxy, a term ed by the floes and fequips, to denote that calmness of mind which feureus us from all emotions arising from vanity and fell-conceit.

\subsection*{ATARNEA}
ATARNEA, an ancient town of Myfn, flatuated between Akrimytition and Plans, remarkable for the marriage of Arribite with the filter or consecutive of the tyrant Hermias; also for the dotage of that phi-lolopier.

\subsection*{ATAXY}
ATAXY, in a general fenfe, the want of order: With phyticians, it signifies irregularity of crises and paroxysites of fevers.

\subsection*{AGCHE}
AGCHE, in commerce, a small liver coin ufed in Turkey, and worth only one third of the English penny.

\subsection*{ATCHIEVEMENT}
ATCHIEVEMENT, in heralyly denotes the arms of a person, or family, together with all the exterior ornament of the fluid; as helmes, mantle, credit, fcrolls, and motto, together with fuch quarterings as may have been acquired by aliiances, all marhalled in order.

\subsection*{ATHCIEVE}
ATHCIEVE. This item is derived from the French \emph{abarne}, i.c. to fish or make an end of his figurities, in heredity, to prefer green schemes or ex-ploids.

\subsection*{ATE}
ATE, the goddels of mithief; in the Pagan tho-ology. She was daughter of Jupiter, and caft down from heaven at the birth of Hercules. For Juno ha-ving decided Jupiter, in scaling Enulflets to be born before Hercules, Jupiter expelled his retirement on Ates, as the author of that miichief; and threw her headlong from heaven to earth, featuring the floodal never return children again (Honeri L six. 125.) The name of the goddels comes from stars, men, to \emph{hart}. Her being the daughter of Jupiter means, according to mythologiils, has no evil happens to us but by the permitief providence, and her healthment to earth denotes the terrible effects of divine jul-tice among men.

\subsection*{ATEGUA}
ATEGUA, or ATEGUA, an ancient town of Spain, placed by fame in the road from Antiquars, now Antquers, a Filipin, or Seville, or others near Aleala Reil; which left in the more probable fi-tuation, because the flumen fallum, now the Salado, was in its neighborhood. Now \emph{Teleda Prins, or Yrtede.}

\subsection*{ATELLA}
ATELLA, an ancient town of Gargania in Italy, between Capua and Nespolis. From this town the \emph{Atlana fabule, or Atelani Inda}, took their name. There were alto called \emph{Ofis}, from their invasion, in whole territory Ashala by. They were generally a species of force, interlarded with much ribalry and buffoonery and foundoreste were ecordix or intestinalis preferred between the each of other plays. The authors thank these forces were not reckoned among the common players, nor deemed infamous , but retained the rights of their tribe, and might be listed for folders, the privilege only of free men. The ruins of this town are full to be feen about 11 miles from the modern Averfa, which was built out of its materials.

\subsection*{A TEMO GUISTO}
A TEMO GUISTO, in multic, figurities to flag or play in an equal, true, and just time.

\subsection*{ATERNUM}
ATERNUM, s town of Lucerna in Italy, now \emph{Alteris}, (Chorusas )). Alfo a town in the territory of the Piceni, now \emph{Pofcara}, a port-town of Naples, fitu-ated and the infinite. E. Long 15.25. V.Lav. 42.50.

\subsection*{ATESTE}
ATESTE, a town in the territory of Veneto in Italy, now called \emph{Rje}. E. Long 12.6. N. Lat. 45.25.

\subsection*{ATHAMADULET}
ATHAMADULET, the prime minister of the Perfin empire, as the grand victor is of the Turkish empire. He is great chancellor of the kingdom, pre-fident of the conced, imperirendment of the finances, and is charged with all foreign affairs.

\subsection*{THAMANTA}
THAMANTA, STORKS, a genus of the di-nynia order, belonging to the pantandria claf of plants. Of this genus Limone enumerated nine species; but none of them most particular notice except there-teals, otherwise called \emph{dance creticus}. This is an unbeliferous plant growing wild in the Levant and the warranza part of Europe. The leaves are irger-larly dipoled, and formed like thole of fence]. The flower-flake ribs about two feet high, finding out many branches, terminated at the top by compound number, compelled of new zo small ones. There have white flowers with the petals, which are focected by sibling, hairy, channelled fruit, divided into two parts, containing one shloog hairy feed. The feeds have a warm biting tails, with an agreeable aromatic fmell. They are kept in the shape as a medium, are corn-native, and laid to be diverted, but are little tied in practice. The plant may be propagated from feeds, which should be flown on an open bed of light dry ground, the following autumn the plants should be taken up, and planted at about a foot dilness in a bed of light sandy earth, where the roots will continue fo-veral years.

\subsection*{ATHANASIA}
ATHANASIA, GOLDLOCKS, a genus of the polygania aqualia order, belonging to the fyngeneria class of plants. Of this genus, Dr.Limous mentions 11 \emph{Species} of which the followingic are by Mr Millerreckoned worthy of a place in those gardens where there are convenience for preferring tender plants, as some of the species will outlive the waters of this country. 1. The dentants grows naturally at the Cape of God Hope. It rises three or four feet high, finding out many fine branches garnished with pretty long narrow leaves, having form resemblance to thole of buckthown plannian. The branches are generally terminated by yellow flowers formed into a kind of tunnel. Thefe flowers appear early in summer, and the feeds ripen in autumn. 3. The trifectate is alio a native of Africa, as are the three following sorts. It is a damby plant, about the height of the former, fending out several weak shoots, garnished with awl-flagged until leaves let close to them, of a pale liberty onborn. The flowers are yellow, and terminate the branches in a corymbus. They appear in August, and the black ripen in Otolom. 3. The ethuthella rides with a strong flambaby halk to the heights of fic or eight feet, fending out many ligneous branches garnished with long narrow leaves terminating in four or five pairs like the hole of fampiro. The flowers are like those of the former, and appear in July and Au-guilt, being feduced by feeds which ripen in Otolber. 3. The pudeless hath frong woody flakes four or five feet high, covered with a woolly bark, as are alfo the branches, which are closely garnished with entire, fearshaped, woolly leaves, and are terminated by yel-low flowers hanging hang food,falls, and are sometimes fucceeded by feeds that ripen in the autumn. 5. The animals is a low annual plants, fellow rings above a foot, high, and feeding forth two or three lender branches garnished with wing-pointed leaves : the flalks are ter-minated
\end{multicols}
\section*{Source record}
\begin{itemize}
\item Page ID: \texttt{britannica:309f765b1ecc83e0a7f7762e44b4a927a7555443:train:2171}
\item Scan: \texttt{page-2171-f8869fa95f19.jpg}, 600 x 832 px
\item Image SHA-256: \texttt{\small f8869fa95f190d0e556bf8d2397c0c36da436ea1b23405995f0aac0f398f26d5}
\item Extraction: \texttt{nvidia/nemotron-3.5-parse}, raw reply \texttt{artifacts/extract/raw/2171.json}
\item Text is Parse OCR output, normalized (long s as s, line-end hyphens joined). It is not corrected by hand.
\end{itemize}
\end{document}
